% ------------------------------------------------------------------------------
\section{Limitaciones conocidas y pendientes}
\label{sec:lim-limitaciones}
% ------------------------------------------------------------------------------

Este capítulo reúne, sin adornos, lo que ConcretoLab Desktop \versionapp{} no hace, lo que hace solo en parte y las
inconsistencias encontradas al revisar el código para escribir esta documentación. Sirve para decidir qué se puede
prometer a un cliente y por dónde seguir.

\begin{advertencia}[No está validado para obras]
Ningún resultado de ConcretoLab sustituye el criterio ni la firma de un ingeniero civil colegiado. El análisis sísmico
es \textbf{preliminar} (anteproyecto y enseñanza): el resumen dice «CUMPLE (preliminar)», el PDF advierte que no es una
memoria de cálculo y el diálogo «Acerca de» repite el aviso. Las pruebas automáticas verifican la programación contra
valores conocidos; no son una validación profesional de los métodos ni de su aplicación a un edificio real.
\end{advertencia}

\subsection{Análisis sísmico}

\begin{itemize}
	\item \textbf{Solo pórticos.} El motor (\ruta{FernandoHuancaAnalisisSismico}) modela columnas y vigas como barras
	      con base empotrada y un diafragma rígido por piso. Los sistemas \ruta{DUAL}, \ruta{MUROS},
	      \ruta{MUROS\_DUCTILIDAD\_LIMITADA} y \ruta{ALBANILERIA} de \ruta{SistemaEstructural} están marcados «Aún no
	      modelable en la aplicación».
	\item \textbf{Los muros no entran en la rigidez.} Un elemento \ruta{MURO} se dibuja y se metra, pero el análisis no
	      incluye su rigidez ni su aporte al cortante; el resultado lo avisa.
	\item \textbf{Losas y zapatas.} Las losas solo se usan para el metrado (el piso se supone diafragma rígido) y las
	      zapatas no se modelan (base empotrada). No hay zonas rígidas en los nudos, interacción suelo-estructura ni diseño
	      de elementos.
	\item \textbf{Tabla 9 y aislamiento.} En categoría A (zonas 2 a 4) y A1 (todas) la E.030 no permite pórticos como
	      único sistema, y para A1 en zonas 3 y 4 exige aislamiento sísmico, que la aplicación no modela. La pantalla lo
	      avisa, pero calcula igual a modo referencial.
	\item \textbf{Cargas.} Las cargas vivas de «Industria o taller común», «Depósito o almacén» y «Edificación esencial»
	      no figuran literalmente en la E.020 (\ruta{verificado = false} en \ruta{UsoEdificacion}); la carga muerta
	      adicional es siempre un valor típico sugerido, no de la norma.
	\item \textbf{Zonificación.} La ubicación por distrito solo existe para el texto 2026 (Anexo II); con la versión 2018
	      la zona se elige a mano.
	\item \textbf{Datos de entrada.} Los parámetros del análisis no se guardan en la base: se ingresan en cada sesión.
	\item \textbf{Recomendación.} «Cómo corregir» solo propone agrandar columnas, en pasos de 5~cm y hasta +60~cm;
	      peraltar vigas o agregar placas queda a criterio del ingeniero.
\end{itemize}

\subsection{Análisis tiempo-historia}

\begin{itemize}
	\item Es \textbf{lineal elástico}, con base empotrada, sin interacción suelo-estructura y con amortiguamiento de
	      Rayleigh del 5~\% ajustado en T1 y 0,2~T1 (entre ambos queda algo por debajo). Las derivas son elásticas.
	\item Es una verificación complementaria y conservadora (R~=~1); \textbf{no sustituye} el análisis modal espectral.
	\item \textbf{No incluye acelerogramas reales}: genera siete sintéticos espectro-compatibles o importa registros del
	      usuario (texto, hasta 10~MB). Con registros que no son espectro-compatibles, el criterio del art. 47.6 (90~\%
	      del espectro en cada registro) se informa como aviso y no bloquea el resultado.
	\item El número mínimo de conjuntos (7 en 2026, 3 en 2018) se comprueba y se avisa si no se alcanza.
\end{itemize}

\subsection{Intercambio CAD y BIM}

\begin{itemize}
	\item \textbf{ETABS (\ruta{.e2k}) es experimental}: nunca se verificó importándolo en ETABS. Los pisos salen de las
	      cotas de vigas y losas y las zapatas no se exportan. El botón, el diálogo y el propio archivo lo indican.
	\item \textbf{IFC2X3} solo se verificó con comprobaciones estructurales del archivo (\ruta{ExportacionIfcEtabsTest}:
	      entidades, referencias resueltas, GUID reproducibles), no abriéndolo en un visor IFC. Todos los elementos
	      cuelgan de un único nivel en la cota 0 y las vigas y losas se extruyen en vertical.
	\item \textbf{DWG no se lee de forma nativa.} Se convierte a DXF con el programa externo \ruta{dwgread.exe} (GNU
	      LibreDWG), que \textbf{no se distribuye} con el instalador: el usuario debe conseguirlo y seleccionarlo o
	      colocarlo en una carpeta \ruta{tools}. Sin él, hay que guardar el plano como DXF ASCII desde AutoCAD.
	\item \textbf{RVT no se lee}: la aplicación explica cómo exportar a DXF desde Revit.
	\item La importación DXF reproduce plantas de pórticos rectangulares con la misma distribución en cada piso. De un
	      bloque solo se toma el punto de inserción, y los textos, cotas, sombreados y el espacio papel se ignoran (el
	      diagnóstico lo muestra).
\end{itemize}

\subsection{Metrado y concreto}

\begin{itemize}
	\item \textbf{Las cantidades de acero son referenciales.} Se estiman con cuantías por tipo (por defecto 60, 150,
	      120, 70 y 80~kg/m$^3$ para zapata, columna, viga, muro y losa), editables en «Cuantías de acero\ldots». No son un
	      metrado de obra.
	\item El \textbf{encofrado} es la superficie de contacto con el molde y no descuenta los contactos entre elementos.
	\item El \textbf{diseño de mezclas} ACI~211.1 cubre concreto de peso normal, f'cr hasta 45~MPa (35 con aire
	      incorporado) y sin aditivos reductores de agua. Los límites de durabilidad se tomaron de ACI~318 / E.060 cap.~4
	      y deben contrastarse con la edición vigente. El resultado es una dosificación de partida que exige mezclas de
	      prueba.
	\item La \textbf{evaluación de resistencia} (E.060, art. 5.6) supone que cada «ensayo» es el promedio de al menos
	      dos probetas; con probetas sueltas el resultado es orientativo.
\end{itemize}

\subsection{Datos}

\begin{itemize}
	\item El esquema \textbf{no tiene claves foráneas}; la integridad depende de los servicios. La entidad
	      \ruta{VinculoGeometrico} declara un \ruta{@UniqueConstraint} (columna, viga) que la migración no crea.
	\item Desde la interfaz \textbf{solo se pueden eliminar proyectos y elementos}: no hay botón para borrar materiales,
	      mezclas, ensayos ni tamizados.
	\item \ruta{ExpedienteEstructural} no tiene pantalla: solo se guarda con \ruta{ExpedienteService} (usado en las
	      pruebas) y viaja en la exportación de proyectos.
	\item La biblioteca de materiales y los tamizados \textbf{no viajan} en el archivo \ruta{.concretolab}.
	\item La \textbf{bitácora} no registra las ediciones de elementos, mezclas o materiales, la eliminación de copias, la
	      restauración ya aplicada ni la conversión de una base con contraseña.
	\item El mensaje de error de la importación para un tipo de elemento inválido dice «use ZAPATA, COLUMNA o VIGA»,
	      aunque también se aceptan \ruta{LOSA} y \ruta{MURO}.
\end{itemize}

\subsection{Distribución}

\begin{itemize}
	\item \textbf{El instalador no está firmado.} Windows SmartScreen muestra «Windows protegió su PC» y el usuario debe
	      pulsar «Más información» y «Ejecutar de todas formas». Hace falta un certificado de firma de código.
	\item Solo hay distribución para \textbf{Windows}. El JAR que se compila incluye las bibliotecas JavaFX de la
	      plataforma donde se construyó.
	\item \ruta{main.wxs}, \ruta{ShortcutPromptDlg.wxs} e \ruta{InstallDirNotEmptyDlg.wxs} son copias de las plantillas
	      del \ruta{jpackage} del JDK~17, y el instalador depende de WiX~3.x. Cambiar de JDK exige revisarlas.
	\item Si una dependencia nueva necesita un módulo del JDK que no está en la lista de \ruta{jlink}, la versión
	      instalada fallará aunque \ruta{java -jar} funcione: \ruta{CREAR\_DISTRIBUCION\_WINDOWS.ps1 -Verificar} y
	      \ruta{PROBAR\_INSTALADOR\_WINDOWS.ps1} sirven para detectarlo.
\end{itemize}

\subsection{Pruebas}

\begin{itemize}
	\item \ruta{ui/VisorMovimientoTest} no se omite sin pantalla: falla, a diferencia de las otras pruebas gráficas.
	\item Las pruebas con Spring no fijan \ruta{concretolab.datos}: al refrescar «Datos y copias» y el resumen de metrado,
	      \ruta{VistasFxmlTest} \textbf{lee} (sin escribir) la carpeta de datos por defecto.
	\item \ruta{EsquemaMigracionesTest} comprueba en el historial de Flyway las versiones 1 y 2, no la 3; V3 solo se
	      ejercita de forma indirecta en \ruta{CompatibilidadMezclaTest}.
\end{itemize}

\subsection{Código sin uso o heredado}

Estas piezas compilan y algunas tienen pruebas, pero ninguna pantalla las usa. Conviene decidir si se conectan o se
retiran:

\begin{longtable}{L{5.2cm}L{9.6cm}}
	\caption{Código sin uso en la aplicación.} \label{tab:lim-sin-uso} \\
	\hline
	\textbf{Elemento} & \textbf{Situación} \\ \hline
	\endfirsthead
	\hline
	\textbf{Elemento} & \textbf{Situación} \\ \hline
	\endhead
	\hline
	\endfoot
	\ruta{AnalisisPortico}, \ruta{EdificioCortante}, \ruta{HistoriaPiso}, \ruta{OsciladorSismico} & Simulación de
	pórtico inicial. Solo los usan sus pruebas (\ruta{HistoriaPiso} lo usa \ruta{EdificioCortante}). \\
	\ruta{DiagnosticoSismico} & Ninguna clase ni prueba lo usa. \\
	\ruta{ui.\allowbreak{}Autocompletar} & Componente de autocompletado sin uso. \\
	\ruta{exception.\allowbreak{}NegocioException} & Tres controladores la capturan, pero nada la lanza. \\
	\ruta{Exportador.Formato.EXCEL} & El \ruta{Exportador} declara el formato \ruta{.xlsx} y su Javadoc habla de Excel,
	pero no hay generador de Excel ni dependencia para ello (Apache POI). \\
	\ruta{BaseCifrada.\allowbreak{}erickDelgadoCifrarExistente}, \ruta{BaseCifrada.\allowbreak{}erickDelgadoCambiarClave} & De cuando la base se
	cifraba; solo las usa \ruta{CifradoBaseTest}. \\
\end{longtable}

\subsection{Documentación desactualizada o inconsistente}

\begin{itemize}
	\item \textbf{La migración V3.} El README dice que el esquema lo gestiona Flyway en
	      \ruta{src/main/resources/db/migration} y enumera «V1 tablas, V2 bitácora, V3 columnas faltantes». En esa carpeta
	      solo hay V1 y V2; \textbf{V3 sí existe}, pero es una migración Java en
	      \ruta{src/main/java/db/migration/V3\_\_Completar\_columnas.java}. El comentario de \ruta{application.properties}
	      («se completa con V1 y V2») no la menciona.
	\item \textbf{Párrafos antiguos del README.} Algunos párrafos iniciales contradicen las secciones posteriores: dicen
	      que la suma de volúmenes es bruta y no descuenta intersecciones (hoy existe el metrado neto), que el módulo de
	      mezclas no diseña una mezcla ACI ni dictamina conformidad (hoy existen «Diseñar mezcla» y la evaluación E.060)
	      y que la granulometría no calcula el módulo de finura (hoy lo calcula y verifica el huso).
	\item \ruta{doc/REFACTORIZACION\_PREFIJOS.md} nombra la clase \ruta{ClaveBaseService}, que ya no existe.
	\item \textbf{«Tú» en lugar de «usted».} La convención es tratar de «usted», pero varios textos usan «tú»: casi toda
	      la pantalla «Datos y copias» (vista y controlador), las indicaciones de \ruta{inicio-view},
	      \ruta{granulometria-view}, \ruta{mezclas-view} y la plantilla de \ruta{modelo-view}, el diálogo de ejes y
	      niveles, el panel «Mover en 3D», el asistente de importación CAD y algunos mensajes de \ruta{ModeloController}.
	      El README también usa «tú».
\end{itemize}

\subsection{Marcadores en el código}

Para completar esta lista se buscó en \ruta{src/main} \ruta{TODO}, \ruta{FIXME}, «pendiente», «experimental», «no
verificado», «preliminar» y «referencial». No hay ningún \ruta{TODO} ni \ruta{FIXME}. Los demás marcadores coinciden con
lo descrito arriba: «experimental» solo aparece en la exportación ETABS; «NO VERIFICADO» en las cargas de
\ruta{UsoEdificacion}; «Aún no modelable» en \ruta{SistemaEstructural}; «preliminar» y «referencial» en los avisos del
análisis sísmico, el metrado y la guía sísmica.
